% Technical Design Document -- CampusGPT
% Build (from this directory):  tectonic -X compile campus-gpt-technical-design.tex
\documentclass[11pt,a4paper,oneside]{report}
\newcommand{\DocShort}{Technical Design}
% Shared preamble for CampusGPT documents (SRS, Technical Design).
% Usage, after \documentclass:
%   \newcommand{\DocShort}{SRS}        % left header text
%   % Shared preamble for CampusGPT documents (SRS, Technical Design).
% Usage, after \documentclass:
%   \newcommand{\DocShort}{SRS}        % left header text
%   % Shared preamble for CampusGPT documents (SRS, Technical Design).
% Usage, after \documentclass:
%   \newcommand{\DocShort}{SRS}        % left header text
%   \input{../common/preamble}

% ---------------------------------------------------------------------------
% Team details
% ---------------------------------------------------------------------------
\newcommand{\Fill}[1]{\textcolor{red}{[#1]}}
\newcommand{\MemberOne}{Ankur Karn & 500118986 & 11}
\newcommand{\MemberTwo}{Ashmita Sharma & 500122422 & 11}
\newcommand{\MemberThree}{Ananya Singh & 500123837 & 11}
\newcommand{\MemberFour}{Ayushi Mishra & 500121986 & 11}
\newcommand{\Specialization}{Cyber Security}
\newcommand{\Mentor}{Mr.~Yogesh Rathia}
\newcommand{\School}{School of Computer Science}
\newcommand{\University}{UNIVERSITY OF PETROLEUM \& ENERGY STUDIES}
\newcommand{\UniversityAddress}{Dehradun -- 248007, Uttarakhand}

% ---------------------------------------------------------------------------
% Packages
% ---------------------------------------------------------------------------
\usepackage[margin=2.5cm,headheight=14pt]{geometry}
\usepackage[T1]{fontenc}
\usepackage{lmodern}
\usepackage{microtype}
\usepackage[table]{xcolor}
\usepackage{booktabs}
\usepackage{tabularx}
\usepackage{array}
\usepackage{longtable}
\usepackage{enumitem}
\usepackage{amsmath}
\usepackage{fancyhdr}
\usepackage{titlesec}
\usepackage{tikz}
\usepackage{pgfgantt}
\usepackage[most]{tcolorbox}
\usepackage{caption}
\usepackage{float}
\usepackage{listings}
\usepackage[hidelinks]{hyperref}
\usepackage{xurl}

\input{../common/diagram-preamble}
\captionsetup[table]{position=above,skip=4pt}
\setlength{\parindent}{0pt}
\setlength{\parskip}{6pt}
\renewcommand{\arraystretch}{1.15}
\newcolumntype{Y}{>{\raggedright\arraybackslash}X}

% Chapter headings: "CHAPTER 1" / "APPENDIX A" above the title
\titleformat{\chapter}[display]
  {\normalfont\bfseries}{\large\MakeUppercase{\chaptertitlename}~\thechapter}{4pt}{\Huge}
\titlespacing*{\chapter}{0pt}{0pt}{24pt}

% Header / footer
\pagestyle{fancy}
\fancyhf{}
\fancyhead[L]{\itshape\DocShort}
\fancyhead[R]{\itshape CampusGPT}
\fancyfoot[C]{Page \thepage}
\renewcommand{\headrulewidth}{0.4pt}
\fancypagestyle{plain}{%
  \fancyhf{}\fancyhead[L]{\itshape\DocShort}\fancyhead[R]{\itshape CampusGPT}%
  \fancyfoot[C]{Page \thepage}\renewcommand{\headrulewidth}{0.4pt}}

% Algorithm box
\newtcolorbox{algobox}[1]{enhanced,breakable,colback=white,colframe=algogreen,
  coltitle=white,fonttitle=\bfseries\sffamily,title={#1},arc=2mm,boxrule=0.6pt,
  left=6pt,right=6pt,top=4pt,bottom=4pt}

% Code listings
\lstset{basicstyle=\ttfamily\footnotesize, columns=fullflexible, keepspaces=true,
  frame=single, rulecolor=\color{gray!50}, backgroundcolor=\color{gray!4},
  keywordstyle=\color{algogreen}\bfseries, commentstyle=\color{gray}\itshape,
  stringstyle=\color{red!50!black}, showstringspaces=false, breaklines=true,
  xleftmargin=4pt, xrightmargin=4pt, aboveskip=6pt, belowskip=6pt,
  captionpos=t}
\renewcommand{\lstlistingname}{Listing}
\renewcommand{\lstlistlistingname}{List of Listings}

% ---------------------------------------------------------------------------
% Title page: \MakeTitlePage{line 1}{line 2 (may be empty)}{version}
% ---------------------------------------------------------------------------
\newcommand{\MakeTitlePage}[3]{%
\hypersetup{pageanchor=false}%
\begin{titlepage}
\thispagestyle{empty}
\centering
\vspace*{0.5cm}
{\Huge\bfseries #1\ifx&#2&\else\\[4pt] #2\fi\par}
\vspace{1.2cm}
{\LARGE\bfseries CampusGPT\par}
\vspace{0.3cm}
{\large\itshape An AI-Powered University Intelligence and Student Assistance Platform\par}
\vspace{1cm}
Submitted in partial fulfilment of the requirements for the award of the degree of\par
\vspace{0.3cm}
{\Large\bfseries Bachelor of Technology\par}
\vspace{0.15cm}
Computer Science and Engineering\par
Specialization: \textbf{\Specialization}\par
\vspace{1cm}
{\large\bfseries Prepared by\par}
\vspace{0.3cm}
\begin{tabular*}{0.8\textwidth}{@{\extracolsep{\fill}}lrr}
\toprule
\textbf{Name} & \textbf{SAP ID} & \textbf{Batch}\\
\midrule
\MemberOne\\
\MemberTwo\\
\MemberThree\\
\MemberFour\\
\bottomrule
\end{tabular*}\par
\vspace{1cm}
Under the guidance of\par
\vspace{0.15cm}
{\large\bfseries \Mentor\par}
\textit{Project Mentor}\par
\vfill
{\bfseries \School\par}
{\bfseries \University\par}
\UniversityAddress\par
\vspace{0.8cm}
{\large\bfseries October 2026 \quad$|$\quad Version #3\par}
\end{titlepage}
\hypersetup{pageanchor=true}}


% ---------------------------------------------------------------------------
% Team details
% ---------------------------------------------------------------------------
\newcommand{\Fill}[1]{\textcolor{red}{[#1]}}
\newcommand{\MemberOne}{Ankur Karn & 500118986 & 11}
\newcommand{\MemberTwo}{Ashmita Sharma & 500122422 & 11}
\newcommand{\MemberThree}{Ananya Singh & 500123837 & 11}
\newcommand{\MemberFour}{Ayushi Mishra & 500121986 & 11}
\newcommand{\Specialization}{Cyber Security}
\newcommand{\Mentor}{Mr.~Yogesh Rathia}
\newcommand{\School}{School of Computer Science}
\newcommand{\University}{UNIVERSITY OF PETROLEUM \& ENERGY STUDIES}
\newcommand{\UniversityAddress}{Dehradun -- 248007, Uttarakhand}

% ---------------------------------------------------------------------------
% Packages
% ---------------------------------------------------------------------------
\usepackage[margin=2.5cm,headheight=14pt]{geometry}
\usepackage[T1]{fontenc}
\usepackage{lmodern}
\usepackage{microtype}
\usepackage[table]{xcolor}
\usepackage{booktabs}
\usepackage{tabularx}
\usepackage{array}
\usepackage{longtable}
\usepackage{enumitem}
\usepackage{amsmath}
\usepackage{fancyhdr}
\usepackage{titlesec}
\usepackage{tikz}
\usepackage{pgfgantt}
\usepackage[most]{tcolorbox}
\usepackage{caption}
\usepackage{float}
\usepackage{listings}
\usepackage[hidelinks]{hyperref}
\usepackage{xurl}

% Settings shared by every diagram (used by the documents and the slide images).
\usetikzlibrary{positioning,arrows.meta,shapes.geometric,fit,calc,backgrounds}
\definecolor{algogreen}{HTML}{1F4D3A}
\definecolor{tierblue}{HTML}{EEF3FA}
\definecolor{tiergreen}{HTML}{EEF6F1}
\definecolor{tieramber}{HTML}{FBF4E6}
\definecolor{tiergray}{HTML}{F2F2F2}
% Numbered channel marker used in diagrams, e.g. \circled{1}
\newcommand{\circled}[1]{\textbf{(#1)}}
% Never hyphenate words inside diagram labels
\tikzset{every node/.append style={execute at begin node={\hyphenpenalty=10000\exhyphenpenalty=10000}}}

\captionsetup[table]{position=above,skip=4pt}
\setlength{\parindent}{0pt}
\setlength{\parskip}{6pt}
\renewcommand{\arraystretch}{1.15}
\newcolumntype{Y}{>{\raggedright\arraybackslash}X}

% Chapter headings: "CHAPTER 1" / "APPENDIX A" above the title
\titleformat{\chapter}[display]
  {\normalfont\bfseries}{\large\MakeUppercase{\chaptertitlename}~\thechapter}{4pt}{\Huge}
\titlespacing*{\chapter}{0pt}{0pt}{24pt}

% Header / footer
\pagestyle{fancy}
\fancyhf{}
\fancyhead[L]{\itshape\DocShort}
\fancyhead[R]{\itshape CampusGPT}
\fancyfoot[C]{Page \thepage}
\renewcommand{\headrulewidth}{0.4pt}
\fancypagestyle{plain}{%
  \fancyhf{}\fancyhead[L]{\itshape\DocShort}\fancyhead[R]{\itshape CampusGPT}%
  \fancyfoot[C]{Page \thepage}\renewcommand{\headrulewidth}{0.4pt}}

% Algorithm box
\newtcolorbox{algobox}[1]{enhanced,breakable,colback=white,colframe=algogreen,
  coltitle=white,fonttitle=\bfseries\sffamily,title={#1},arc=2mm,boxrule=0.6pt,
  left=6pt,right=6pt,top=4pt,bottom=4pt}

% Code listings
\lstset{basicstyle=\ttfamily\footnotesize, columns=fullflexible, keepspaces=true,
  frame=single, rulecolor=\color{gray!50}, backgroundcolor=\color{gray!4},
  keywordstyle=\color{algogreen}\bfseries, commentstyle=\color{gray}\itshape,
  stringstyle=\color{red!50!black}, showstringspaces=false, breaklines=true,
  xleftmargin=4pt, xrightmargin=4pt, aboveskip=6pt, belowskip=6pt,
  captionpos=t}
\renewcommand{\lstlistingname}{Listing}
\renewcommand{\lstlistlistingname}{List of Listings}

% ---------------------------------------------------------------------------
% Title page: \MakeTitlePage{line 1}{line 2 (may be empty)}{version}
% ---------------------------------------------------------------------------
\newcommand{\MakeTitlePage}[3]{%
\hypersetup{pageanchor=false}%
\begin{titlepage}
\thispagestyle{empty}
\centering
\vspace*{0.5cm}
{\Huge\bfseries #1\ifx&#2&\else\\[4pt] #2\fi\par}
\vspace{1.2cm}
{\LARGE\bfseries CampusGPT\par}
\vspace{0.3cm}
{\large\itshape An AI-Powered University Intelligence and Student Assistance Platform\par}
\vspace{1cm}
Submitted in partial fulfilment of the requirements for the award of the degree of\par
\vspace{0.3cm}
{\Large\bfseries Bachelor of Technology\par}
\vspace{0.15cm}
Computer Science and Engineering\par
Specialization: \textbf{\Specialization}\par
\vspace{1cm}
{\large\bfseries Prepared by\par}
\vspace{0.3cm}
\begin{tabular*}{0.8\textwidth}{@{\extracolsep{\fill}}lrr}
\toprule
\textbf{Name} & \textbf{SAP ID} & \textbf{Batch}\\
\midrule
\MemberOne\\
\MemberTwo\\
\MemberThree\\
\MemberFour\\
\bottomrule
\end{tabular*}\par
\vspace{1cm}
Under the guidance of\par
\vspace{0.15cm}
{\large\bfseries \Mentor\par}
\textit{Project Mentor}\par
\vfill
{\bfseries \School\par}
{\bfseries \University\par}
\UniversityAddress\par
\vspace{0.8cm}
{\large\bfseries October 2026 \quad$|$\quad Version #3\par}
\end{titlepage}
\hypersetup{pageanchor=true}}


% ---------------------------------------------------------------------------
% Team details
% ---------------------------------------------------------------------------
\newcommand{\Fill}[1]{\textcolor{red}{[#1]}}
\newcommand{\MemberOne}{Ankur Karn & 500118986 & 11}
\newcommand{\MemberTwo}{Ashmita Sharma & 500122422 & 11}
\newcommand{\MemberThree}{Ananya Singh & 500123837 & 11}
\newcommand{\MemberFour}{Ayushi Mishra & 500121986 & 11}
\newcommand{\Specialization}{Cyber Security}
\newcommand{\Mentor}{Mr.~Yogesh Rathia}
\newcommand{\School}{School of Computer Science}
\newcommand{\University}{UNIVERSITY OF PETROLEUM \& ENERGY STUDIES}
\newcommand{\UniversityAddress}{Dehradun -- 248007, Uttarakhand}

% ---------------------------------------------------------------------------
% Packages
% ---------------------------------------------------------------------------
\usepackage[margin=2.5cm,headheight=14pt]{geometry}
\usepackage[T1]{fontenc}
\usepackage{lmodern}
\usepackage{microtype}
\usepackage[table]{xcolor}
\usepackage{booktabs}
\usepackage{tabularx}
\usepackage{array}
\usepackage{longtable}
\usepackage{enumitem}
\usepackage{amsmath}
\usepackage{fancyhdr}
\usepackage{titlesec}
\usepackage{tikz}
\usepackage{pgfgantt}
\usepackage[most]{tcolorbox}
\usepackage{caption}
\usepackage{float}
\usepackage{listings}
\usepackage[hidelinks]{hyperref}
\usepackage{xurl}

% Settings shared by every diagram (used by the documents and the slide images).
\usetikzlibrary{positioning,arrows.meta,shapes.geometric,fit,calc,backgrounds}
\definecolor{algogreen}{HTML}{1F4D3A}
\definecolor{tierblue}{HTML}{EEF3FA}
\definecolor{tiergreen}{HTML}{EEF6F1}
\definecolor{tieramber}{HTML}{FBF4E6}
\definecolor{tiergray}{HTML}{F2F2F2}
% Numbered channel marker used in diagrams, e.g. \circled{1}
\newcommand{\circled}[1]{\textbf{(#1)}}
% Never hyphenate words inside diagram labels
\tikzset{every node/.append style={execute at begin node={\hyphenpenalty=10000\exhyphenpenalty=10000}}}

\captionsetup[table]{position=above,skip=4pt}
\setlength{\parindent}{0pt}
\setlength{\parskip}{6pt}
\renewcommand{\arraystretch}{1.15}
\newcolumntype{Y}{>{\raggedright\arraybackslash}X}

% Chapter headings: "CHAPTER 1" / "APPENDIX A" above the title
\titleformat{\chapter}[display]
  {\normalfont\bfseries}{\large\MakeUppercase{\chaptertitlename}~\thechapter}{4pt}{\Huge}
\titlespacing*{\chapter}{0pt}{0pt}{24pt}

% Header / footer
\pagestyle{fancy}
\fancyhf{}
\fancyhead[L]{\itshape\DocShort}
\fancyhead[R]{\itshape CampusGPT}
\fancyfoot[C]{Page \thepage}
\renewcommand{\headrulewidth}{0.4pt}
\fancypagestyle{plain}{%
  \fancyhf{}\fancyhead[L]{\itshape\DocShort}\fancyhead[R]{\itshape CampusGPT}%
  \fancyfoot[C]{Page \thepage}\renewcommand{\headrulewidth}{0.4pt}}

% Algorithm box
\newtcolorbox{algobox}[1]{enhanced,breakable,colback=white,colframe=algogreen,
  coltitle=white,fonttitle=\bfseries\sffamily,title={#1},arc=2mm,boxrule=0.6pt,
  left=6pt,right=6pt,top=4pt,bottom=4pt}

% Code listings
\lstset{basicstyle=\ttfamily\footnotesize, columns=fullflexible, keepspaces=true,
  frame=single, rulecolor=\color{gray!50}, backgroundcolor=\color{gray!4},
  keywordstyle=\color{algogreen}\bfseries, commentstyle=\color{gray}\itshape,
  stringstyle=\color{red!50!black}, showstringspaces=false, breaklines=true,
  xleftmargin=4pt, xrightmargin=4pt, aboveskip=6pt, belowskip=6pt,
  captionpos=t}
\renewcommand{\lstlistingname}{Listing}
\renewcommand{\lstlistlistingname}{List of Listings}

% ---------------------------------------------------------------------------
% Title page: \MakeTitlePage{line 1}{line 2 (may be empty)}{version}
% ---------------------------------------------------------------------------
\newcommand{\MakeTitlePage}[3]{%
\hypersetup{pageanchor=false}%
\begin{titlepage}
\thispagestyle{empty}
\centering
\vspace*{0.5cm}
{\Huge\bfseries #1\ifx&#2&\else\\[4pt] #2\fi\par}
\vspace{1.2cm}
{\LARGE\bfseries CampusGPT\par}
\vspace{0.3cm}
{\large\itshape An AI-Powered University Intelligence and Student Assistance Platform\par}
\vspace{1cm}
Submitted in partial fulfilment of the requirements for the award of the degree of\par
\vspace{0.3cm}
{\Large\bfseries Bachelor of Technology\par}
\vspace{0.15cm}
Computer Science and Engineering\par
Specialization: \textbf{\Specialization}\par
\vspace{1cm}
{\large\bfseries Prepared by\par}
\vspace{0.3cm}
\begin{tabular*}{0.8\textwidth}{@{\extracolsep{\fill}}lrr}
\toprule
\textbf{Name} & \textbf{SAP ID} & \textbf{Batch}\\
\midrule
\MemberOne\\
\MemberTwo\\
\MemberThree\\
\MemberFour\\
\bottomrule
\end{tabular*}\par
\vspace{1cm}
Under the guidance of\par
\vspace{0.15cm}
{\large\bfseries \Mentor\par}
\textit{Project Mentor}\par
\vfill
{\bfseries \School\par}
{\bfseries \University\par}
\UniversityAddress\par
\vspace{0.8cm}
{\large\bfseries October 2026 \quad$|$\quad Version #3\par}
\end{titlepage}
\hypersetup{pageanchor=true}}


\begin{document}

\MakeTitlePage{Technical Design}{Document}{1.0}

\pagenumbering{roman}
\setcounter{page}{1}
\tableofcontents
\listoffigures
\listoftables
\lstlistoflistings

\chapter*{Revision History}
\addcontentsline{toc}{chapter}{Revision History}
\begin{tabularx}{\textwidth}{|c|Y|Y|}
\hline
\textbf{Date} & \textbf{Change} & \textbf{Reason for Changes}\\
\hline
Tue, 13 Oct 2026 & Version 1.0. Technical diagrams, inter-process communication design,
programming concepts and library selection. & First submission for evaluation, alongside
SRS version 1.1.\\
\hline
\end{tabularx}

\chapter*{Document Conventions}
\addcontentsline{toc}{chapter}{Document Conventions}
\begin{itemize}
  \item This document describes \emph{how} CampusGPT is built. \emph{What} it must do is
        specified in the Software Requirements Specification (SRS); requirement IDs such
        as \textbf{FR-08} refer to that document.
  \item Code listings are design sketches in Python and SQL. They show the intended
        structure; identifiers may change during implementation.
  \item Numbered channel markers such as \circled{3} identify the same inter-process
        communication channel in every diagram and table.
  \item All student records and policy documents are synthetic (fictional).
\end{itemize}

% ===========================================================================
\clearpage
\pagenumbering{arabic}
\setcounter{page}{1}
\chapter{Introduction}

\section{Purpose}
CampusGPT is an AI-powered helpdesk that answers university students' questions from two
sources: their own records, read through permission-checked backend tools, and policy
documents, searched with Retrieval-Augmented Generation (RAG)~\cite{lewis2020rag}. This
document presents the technical design: the system's structure and data, the processes
that make up the running system and how they communicate, the programming concepts used to
build it, and the libraries chosen.

\section{Evaluation Components}
Table~\ref{tab:mapping} shows where each evaluated component is addressed.

\begin{table}[H]
\centering
\caption{Evaluation components and where they are addressed}
\label{tab:mapping}
\begin{tabularx}{\textwidth}{@{}lYc@{}}
\toprule
\textbf{Component} & \textbf{Content} & \textbf{Chapter}\\
\midrule
Technical Diagram & System architecture, Level-1 DFD, sequence diagram, ER diagrams,
process view, job state machine & 2, 3\\
Programming Concepts & Object-oriented design, design patterns, asynchronous programming,
generators, validation, transactions, data structures, testing & 4\\
Inter-Process Communication & Seven IPC mechanisms across five processes, with the
document-ingestion pipeline as the worked example & 3\\
Libraries & Each library's purpose, the alternative considered, and the reason for the
choice & 5\\
SRS Document & Submitted separately (version 1.1) & --\\
\bottomrule
\end{tabularx}
\end{table}

\section{Design Principles}
\begin{enumerate}[itemsep=1pt]
  \item \textbf{The language model never touches the database.} It can only request one of
        a fixed set of read-only tools; the backend checks permission and supplies the
        student's identity from the login session.
  \item \textbf{Facts and calculations come from code.} Numbers are read from PostgreSQL and
        computed deterministically; the model selects tools and words the answer.
  \item \textbf{Slow work runs in a separate process.} Document ingestion is handed to a
        worker process so that chat requests are never blocked.
  \item \textbf{One database for everything.} PostgreSQL stores records, vectors, full-text
        indexes and the job queue, and also carries notifications between processes.
  \item \textbf{Replaceable external parts.} The LLM provider and the university data
        source sit behind interfaces and can be swapped by configuration.
\end{enumerate}

% ===========================================================================
\chapter{Technical Diagrams}

\section{System Architecture}
Figure~\ref{fig:arch} shows CampusGPT as four tiers. Grey arrows are function calls inside
one process; black arrows, numbered \circled{1}--\circled{4}, cross a process boundary and
are detailed in Chapter~\ref{ch:ipc}.

\begin{figure}[H]
\centering
\resizebox{\textwidth}{!}{% System architecture diagram (shared by SRS and Technical Design)
\begin{tikzpicture}[
  font=\sffamily\scriptsize,
  comp/.style={draw=black!60, fill=white, rounded corners=2pt, align=center,
               text width=2.75cm, minimum height=0.95cm, inner sep=2pt},
  tier/.style={draw=black!50, rounded corners=4pt, inner sep=6pt},
  tierlabel/.style={font=\sffamily\scriptsize\bfseries, anchor=north west},
  ext/.style={draw=black!70, fill=tiergray, rounded corners=2pt, align=center,
              text width=2.75cm, minimum height=1cm, inner sep=2pt},
  arr/.style={-{Latex[length=2mm]}, thick, draw=black!75},
  biarr/.style={{Latex[length=2mm]}-{Latex[length=2mm]}, thick, draw=black!75},
  lab/.style={font=\sffamily\tiny, fill=white, inner sep=1.5pt, align=center}]

  % ---------------- Client tier ----------------
  \node[comp] (stuui) at (-4.9,10.1) {\textbf{Student Portal}\\Dashboard, Chat};
  \node[comp] (admui) at (-1.4,10.1) {\textbf{Admin Console}\\Uploads, Status};
  \node[comp] (next)  at (2.1,10.1)  {\textbf{Next.js App}\\React + TypeScript};
  \begin{scope}[on background layer]
    \node[tier, fill=tierblue, fit=(stuui)(admui)(next), inner ysep=9pt,
          yshift=3pt] (client) {};
  \end{scope}
  \node[tierlabel] at (client.north west) {Client Tier (Web Browser)};

  % ---------------- External LLM ----------------
  \node[ext] (llm) at (6.05,10.1) {\textbf{LLM Provider}\\Groq, Gemini\\or Ollama};

  % ---------------- Application tier ----------------
  \node[comp] (auth)  at (-4.9,7.6) {\textbf{Auth \& RBAC}\\Session check, roles};
  \node[comp] (api)   at (-1.4,7.6) {\textbf{API Routers}\\REST + SSE};
  \node[comp] (adapt) at (2.1,7.6)  {\textbf{LLM Adapter}\\Provider interface};
  \node[comp] (orch)  at (-4.9,6.1) {\textbf{Chat Orchestrator}\\Tool loop, citations};
  \node[comp] (tools) at (-1.4,6.1) {\textbf{Tool Registry}\\9 read-only tools};
  \node[comp] (retr)  at (2.1,6.1)  {\textbf{Retrieval Service}\\Hybrid search (RRF)};
  \node[comp] (week)  at (-4.9,4.6) {\textbf{Weekly Summary}\\Urgency rules};
  \node[comp] (calc)  at (-1.4,4.6) {\textbf{Calculators}\\Attendance formula};
  \node[comp] (emb)   at (2.1,4.6)  {\textbf{Query Embedder}\\BGE-small (CPU)};
  \begin{scope}[on background layer]
    \node[tier, fill=tiergreen, fit=(auth)(adapt)(week)(emb), inner ysep=9pt,
          yshift=3pt] (app) {};
  \end{scope}
  \node[tierlabel] at (app.north west) {Application Tier (FastAPI process)};

  % ---------------- Worker ----------------
  \node[comp] (listen) at (6.05,7.6) {\textbf{Job Listener}\\LISTEN + claim job};
  \node[comp] (extract) at (6.05,6.1) {\textbf{Extract \& Chunk}\\pypdf, cleaning};
  \node[comp] (wemb) at (6.05,4.6) {\textbf{Chunk Embedder}\\BGE-small (CPU)};
  \begin{scope}[on background layer]
    \node[tier, fill=tieramber, fit=(listen)(wemb), inner ysep=9pt, yshift=3pt]
      (worker) {};
  \end{scope}
  \node[tierlabel] at (worker.north west) {Ingestion Worker};

  % ---------------- Data tier ----------------
  \node[comp] (rec)  at (-4.9,1.7) {\textbf{Academic Records}\\students, courses, fees};
  \node[comp] (vec)  at (-1.4,1.7) {\textbf{Document Chunks}\\pgvector + full-text};
  \node[comp] (jobs) at (2.1,1.7)  {\textbf{Job Queue}\\ingestion\_jobs};
  \node[comp] (logs) at (6.05,1.7) {\textbf{Chat \& Audit Logs}\\Conversations, events};
  \begin{scope}[on background layer]
    \node[tier, fill=tiergray, fit=(rec)(logs), inner ysep=9pt, yshift=3pt]
      (data) {};
  \end{scope}
  \node[tierlabel] at (data.north west) {Data Tier (PostgreSQL 16 + pgvector)};

  % ---------------- Inter-tier connections ----------------
  \draw[biarr] (api.north) -- node[lab] {\circled{1}\,HTTPS: REST (JSON)\\+ SSE answer stream}
    (api.north|-client.south);
  \draw[biarr] (adapt.east) -- ++(0.35,0) |- node[lab, pos=0.25]
    {\circled{4}\,HTTPS\\JSON} (llm.west);
  \draw[biarr] (tools.south|-app.south) -- node[lab]
    {\circled{2}\,SQL over TCP\\(psycopg pool)} (tools.south|-data.north);
  \draw[biarr] (wemb.south|-worker.south) -- node[lab]
    {\circled{3}\,LISTEN/NOTIFY\\+ SQL} (wemb.south|-data.north);

  % ---------------- Intra-process flow ----------------
  \draw[arr, draw=black!45] (api) -- (orch);
  \draw[arr, draw=black!45] (orch) -- (tools);
  \draw[arr, draw=black!45] (tools) -- (retr);
  \draw[arr, draw=black!45] (orch.south) -- (week.north);
  \draw[arr, draw=black!45] (tools) -- (calc);
  \draw[arr, draw=black!45] (retr) -- (emb);
  \draw[arr, draw=black!45] (orch.north east) -- (adapt.south west);
  \draw[arr, draw=black!45] (auth) -- (api);
  \draw[arr, draw=black!45] (listen) -- (extract);
  \draw[arr, draw=black!45] (extract) -- (wemb);
\end{tikzpicture}
}
\caption{System architecture of CampusGPT.}
\label{fig:arch}
\end{figure}

\begin{itemize}[itemsep=1pt]
  \item \textbf{Client tier.} A Next.js application with a student portal (dashboard and
        chat) and an admin console (document upload and status).
  \item \textbf{Application tier.} One FastAPI process. Every request passes the
        authentication and role check first. The chat orchestrator runs the tool loop with
        the LLM, the tool registry executes permitted tools, and the retrieval service
        performs hybrid search using a locally loaded embedding model.
  \item \textbf{Ingestion worker.} A separate Python process that converts uploaded
        documents into searchable chunks, so CPU-heavy embedding never slows down chat.
  \item \textbf{Data tier.} PostgreSQL~16 with the pgvector extension~\cite{pgvector} holds
        academic records, document chunks with vectors, the job queue and logs.
\end{itemize}

\section{Data Flow Diagram (Level 1)}
The Level-0 context diagram appears in the SRS. Figure~\ref{fig:dfd} decomposes the system
into five processes and five data stores, following DeMarco's structured-analysis
notation~\cite{demarco1978}. Every data flow passes through a process; no external entity
reads or writes a data store directly.

\begin{figure}[H]
\centering
\resizebox{\textwidth}{!}{% Level-1 Data Flow Diagram (shared by SRS and Technical Design)
\begin{tikzpicture}[
  font=\sffamily\scriptsize,
  proc/.style={circle, draw=algogreen, thick, fill=tiergreen, align=center,
               text width=1.75cm, inner sep=1pt},
  ent/.style={draw=black!70, fill=tierblue, align=center, text width=2.1cm,
              minimum height=0.9cm},
  store/.style={align=left, text width=2.9cm, minimum height=0.6cm, inner xsep=4pt},
  arr/.style={-{Latex[length=2mm]}, draw=black!75},
  lab/.style={font=\sffamily\tiny, fill=white, inner sep=1pt, align=center}]

  % External entities
  \node[ent] (stu) at (-7.2,3.6) {\textbf{Student}};
  \node[ent] (adm) at (7.4,-3.6) {\textbf{Administrator}};
  \node[ent] (llm) at (-2.6,-1.6) {\textbf{LLM Provider}};

  % Processes
  \node[proc] (p1) at (-2.6,6.4) {\textbf{1.0}\\Authenticate User};
  \node[proc] (p2) at (-2.6,2.4) {\textbf{2.0}\\Orchestrate Answer};
  \node[proc] (p3) at (2.4,4.4)  {\textbf{3.0}\\Execute Student Tools};
  \node[proc] (p4) at (2.4,0.4)  {\textbf{4.0}\\Retrieve Policy Text};
  \node[proc] (p5) at (7.4,0.4)  {\textbf{5.0}\\Ingest Documents};

  % Data stores
  \node[store] (d1) at (2.4,7.3)  {\textbf{D1}\enspace Users \& Sessions};
  \node[store] (d2) at (7.4,4.4)  {\textbf{D2}\enspace Academic Records};
  \node[store] (d3) at (2.4,-2.0) {\textbf{D3}\enspace Documents \& Chunks};
  \node[store] (d4) at (-7.2,0.4) {\textbf{D4}\enspace Conversations \& Audit Log};
  \node[store] (d5) at (7.4,2.7) {\textbf{D5}\enspace Ingestion Job Queue};

  % Student <-> authentication
  \draw[arr] (stu.north) |- node[lab, pos=0.75] {credentials} ([yshift=3pt]p1.west);
  \draw[arr] ([yshift=-6pt]p1.west) -| node[lab, pos=0.3] {session} ([xshift=8pt]stu.north);
  \draw[arr] (p1.east) -- node[lab] {lookup, create session} (d1.west);

  % Student <-> orchestrator
  \draw[arr] ([yshift=4pt]stu.south east) -- node[lab, sloped, above] {question} ([yshift=6pt]p2.west);
  \draw[arr] ([yshift=-6pt]p2.west) -- node[lab, sloped, below] {answer + sources} ([yshift=-6pt]stu.south east);
  \draw[arr] (p1.south) -- node[lab] {user id, role} (p2.north);

  % Orchestrator <-> LLM
  \draw[arr] ([xshift=-8pt]p2.south) -- node[lab, left] {prompt, tool list,\\context} ([xshift=-8pt]llm.north);
  \draw[arr] ([xshift=8pt]llm.north) -- node[lab, right] {tool calls,\\answer} ([xshift=8pt]p2.south);

  % Orchestrator <-> tools / retrieval
  \draw[arr] (p2.40) -- node[lab, sloped, above] {tool + session id} (p3.200);
  \draw[arr] (p3.235) -- node[lab, sloped, below] {records, results} (p2.12);
  \draw[arr] (p2.-12) -- node[lab, sloped, above] {search query} (p4.145);
  \draw[arr] (p4.200) -- node[lab, sloped, below] {ranked chunks} (p2.-40);

  % Tools <-> records
  \draw[arr] (d2.west) -- node[lab] {student's rows only} (p3.east);

  % Orchestrator -> log
  \draw[arr] (p2.230) -- node[lab, sloped, below] {messages, tool events} (d4.east);

  % Retrieval <- documents
  \draw[arr] (d3.north) -- node[lab] {chunks, vectors} (p4.south);

  % Administrator -> ingestion
  \draw[arr] ([xshift=-6pt]adm.north) -- node[lab, left, pos=0.3] {document} ([xshift=-6pt]p5.south);
  \draw[arr] ([xshift=6pt]p5.south) -- node[lab, right, pos=0.7] {status} ([xshift=6pt]adm.north);
  \draw[arr] (p5.south west) |- node[lab, pos=0.75] {chunks + vectors} (d3.east);
  \draw[arr] ([xshift=-10pt]p5.north) -- node[lab, left] {enqueue} ([xshift=-10pt]d5.south);
  \draw[arr] ([xshift=10pt]d5.south) -- node[lab, right] {claim} ([xshift=10pt]p5.north);
  % Data-store outlines (open on the right, as in Gane-Sarson notation)
  \foreach \s in {d1,d2,d3,d4,d5}
    \draw[black!70, thick] (\s.north east) -- (\s.north west) -- (\s.south west)
                           -- (\s.south east);
\end{tikzpicture}
}
\caption{Level-1 Data Flow Diagram of CampusGPT.}
\label{fig:dfd}
\end{figure}

\section{Sequence Diagram: Answering a Question}
Figure~\ref{fig:seq} traces the question ``Can I sit for my DBMS exam?'' from a student
with 41 of 60 classes attended. The orchestrator first obtains the student's attendance,
then the relevant policy text, and only then asks the model for an answer. The student
identifier is added by the orchestrator, never by the model, and every citation is checked
against the chunks that were actually retrieved.

\begin{figure}[H]
\centering
\resizebox{\textwidth}{!}{% Sequence diagram: "Can I sit for my DBMS exam?"
\begin{tikzpicture}[
  font=\sffamily\scriptsize,
  head/.style={draw=black!70, fill=tierblue, align=center, text width=1.9cm,
               minimum height=0.8cm, inner sep=2pt},
  call/.style={-{Latex[length=2mm]}, draw=black!80},
  ret/.style={-{Latex[length=2mm]}, dashed, draw=black!60},
  self/.style={-{Latex[length=2mm]}, draw=black!80, rounded corners=2pt},
  mlab/.style={font=\sffamily\tiny, fill=white, inner sep=1pt, align=center},
  step/.style={font=\sffamily\tiny\bfseries, text=algogreen},
  msglab/.style={font=\sffamily\tiny, inner sep=1pt, above=0.5pt}]

  % Lifelines: x positions
  \def\xB{0}\def\xA{3.1}\def\xO{6.2}\def\xL{9.3}\def\xT{12.4}\def\xD{15.5}
  \def\ybot{-15.3}
  \foreach \x/\name in {\xB/Browser\\(Next.js), \xA/FastAPI\\Routers + Auth,
                        \xO/Chat\\Orchestrator, \xL/LLM Provider\\(via adapter),
                        \xT/Tool Registry\\+ Retrieval, \xD/PostgreSQL\\+ pgvector} {
    \node[head] at (\x,0) {\textbf{\name}};
    \draw[black!40, dashed] (\x,-0.45) -- (\x,\ybot);
  }

  % Messages: \msg{style}{from x}{to x}{y}{label}
  \newcommand{\msg}[5]{\draw[#1] (#2,#4) -- node[msglab] {#5} (#3,#4);}

  \msg{call}{\xB}{\xA}{-1.0}{POST /chat (session cookie)}
  \msg{call}{\xA}{\xD}{-1.6}{look up session}
  \msg{ret}{\xD}{\xA}{-2.2}{user\_id, role = student}
  \msg{call}{\xA}{\xO}{-2.8}{run(question, user)}
  \msg{call}{\xO}{\xL}{-3.4}{question + allowed tools}
  \msg{ret}{\xL}{\xO}{-4.0}{call get\_my\_attendance(course=``DBMS'')}
  \draw[self] (\xO,-4.35) -- ++(0.55,0) -- ++(0,-0.45) node[mlab, right, pos=0.5]
    {allow-list check,\\inject student\_id} -- ++(-0.55,0);
  \msg{call}{\xO}{\xT}{-5.2}{get\_my\_attendance(student\_id, ``DBMS'')}
  \msg{call}{\xT}{\xD}{-5.8}{SELECT \ldots{} WHERE student\_id = \$1}
  \msg{ret}{\xD}{\xT}{-6.4}{41 of 60 classes}
  \msg{ret}{\xT}{\xO}{-7.0}{68.3\%, needs 16 more classes}
  \msg{call}{\xO}{\xL}{-7.6}{tool result}
  \msg{ret}{\xL}{\xO}{-8.2}{call search\_documents(``exam attendance'')}
  \msg{call}{\xO}{\xT}{-8.8}{search\_documents(query)}
  \msg{call}{\xT}{\xD}{-9.4}{vector + full-text search}
  \msg{ret}{\xD}{\xT}{-10.0}{chunks C12, C47}
  \msg{ret}{\xT}{\xO}{-10.6}{ranked chunks (RRF)}
  \msg{call}{\xO}{\xL}{-11.2}{chunks as untrusted context}
  \msg{ret}{\xL}{\xO}{-11.8}{answer citing C12}
  \draw[self] (\xO,-12.15) -- ++(0.55,0) -- ++(0,-0.45) node[mlab, right, pos=0.5]
    {verify citations} -- ++(-0.55,0);
  \msg{call}{\xO}{\xD}{-13.0}{store messages, audit events}
  \msg{ret}{\xO}{\xA}{-13.6}{answer + sources}
  \msg{ret}{\xA}{\xB}{-14.2}{SSE: streamed answer + sources}
\end{tikzpicture}
}
\caption{Sequence diagram for a question that needs both records and policy.}
\label{fig:seq}
\end{figure}

\clearpage
\section{Database Design (ER Diagrams)}
The schema follows the relational model~\cite{codd1970} and is drawn in entity--relationship
form~\cite{chen1976}. Figure~\ref{fig:er-academic} shows the academic records and
Figure~\ref{fig:er-platform} the tables supporting authentication, documents and chat. The
\texttt{users} table appears in both figures.

\begin{figure}[H]
\centering
\resizebox{\textwidth}{!}{% ER diagram: academic records
% Shared styles and macro for ER diagrams.
% \entity{node name}{x,y}{table name}{attribute rows separated by \\}
\tikzset{
  entity/.style={draw=black!70, fill=white, inner sep=0pt, font=\sffamily\tiny},
  rel/.style={draw=black!70, thick},
  card/.style={font=\sffamily\tiny\bfseries, fill=white, inner sep=0.8pt, text=algogreen}}
\providecommand{\entity}[4]{%
  \node[entity] (#1) at (#2) {%
    \renewcommand{\arraystretch}{1.05}\setlength{\tabcolsep}{3pt}%
    \begin{tabular}{l}
      \rowcolor{tierblue}\rule{0pt}{8pt}\textbf{\small #3}\\ \hline
      #4
    \end{tabular}};
  \draw[black!70] (#1.south west) rectangle (#1.north east);}
% \relate{from}{to}{card near from}{card near to}[path options]
\providecommand{\relate}[5][]{%
  \draw[rel, #1] (#2) -- node[card, pos=0.12] {#4} node[card, pos=0.88] {#5} (#3);}
\providecommand{\PK}{\textbf{PK}~}
\providecommand{\FK}{\textit{FK}~}

\begin{tikzpicture}
  % Row 1
  \entity{users}{0,0}{users}{\PK id\\ email (unique)\\ password\_hash\\ role}
  \entity{students}{4,0}{students}{\PK id\\ \FK user\_id (unique)\\ \FK program\_id\\ roll\_no, name\\ current\_semester, section\\ is\_synthetic}
  \entity{programs}{8,0}{programs}{\PK id\\ \FK department\_id\\ name, duration\_semesters}
  \entity{departments}{12,0}{departments}{\PK id\\ code (unique), name}
  \entity{faculty}{16,0}{faculty}{\PK id\\ \FK department\_id\\ name, designation}
  % Row 2
  \entity{invoices}{0,-4.4}{invoices}{\PK id\\ \FK student\_id\\ term, amount\\ due\_date, status}
  \entity{enrollments}{4,-4.4}{enrollments}{\PK id\\ \FK student\_id\\ \FK offering\_id\\ unique(student, offering)}
  \entity{offerings}{8,-4.4}{course\_offerings}{\PK id\\ \FK course\_id\\ \FK faculty\_id\\ term, section}
  \entity{courses}{12,-4.4}{courses}{\PK id\\ \FK department\_id\\ code (unique), title\\ credits, semester}
  % Row 3
  \entity{payments}{0,-8.8}{payments}{\PK id\\ \FK invoice\_id\\ amount, paid\_at}
  \entity{attendance}{4,-8.8}{attendance}{\PK (session\_id,\\ \quad enrollment\_id)\\ present (boolean)}
  \entity{sessions}{8,-8.8}{class\_sessions}{\PK id\\ \FK offering\_id\\ held\_on}
  \entity{exams}{12,-8.8}{exams}{\PK id\\ \FK offering\_id\\ exam\_type, starts\_at\\ room}
  % Row 4
  \entity{submissions}{4,-13.0}{submissions}{\PK id\\ \FK assignment\_id\\ \FK enrollment\_id\\ submitted\_at, marks}
  \entity{assignments}{8,-13.0}{assignments}{\PK id\\ \FK offering\_id\\ title, due\_at\\ max\_marks}

  % Relationships (cardinality: 1 / N)
  \relate{students.west}{users.east}{1}{1}
  \relate{students.east}{programs.west}{N}{1}
  \relate{programs.east}{departments.west}{N}{1}
  \relate{faculty.west}{departments.east}{N}{1}
  \relate{courses.north}{departments.south}{N}{1}
  \relate{offerings.east}{courses.west}{N}{1}
  \relate{offerings.north east}{faculty.south west}{N}{1}
  \relate{enrollments.north}{students.south}{N}{1}
  \relate{enrollments.east}{offerings.west}{N}{1}
  \relate{invoices.north east}{students.south west}{N}{1}
  \relate{payments.north}{invoices.south}{N}{1}
  \relate{attendance.north}{enrollments.south}{N}{1}
  \relate{attendance.east}{sessions.west}{N}{1}
  \relate{sessions.north}{offerings.south}{N}{1}
  \relate{exams.north west}{offerings.south east}{N}{1}
  \relate{submissions.east}{assignments.west}{N}{1}
  \draw[rel] (assignments.east) -- node[card, pos=0.3] {N} ++(0.55,0) |- node[card, pos=0.95] {1}
    ([yshift=-6pt]offerings.east);
  \draw[rel] (submissions.west) -- node[card, pos=0.5] {N} ++(-0.45,0) |- node[card, pos=0.95] {1}
    ([yshift=-8pt]enrollments.west);
\end{tikzpicture}
}
\caption{ER diagram: academic records.}
\label{fig:er-academic}
\end{figure}

\begin{figure}[H]
\centering
\resizebox{0.92\textwidth}{!}{% ER diagram: documents, AI and platform tables
% Shared styles and macro for ER diagrams.
% \entity{node name}{x,y}{table name}{attribute rows separated by \\}
\tikzset{
  entity/.style={draw=black!70, fill=white, inner sep=0pt, font=\sffamily\tiny},
  rel/.style={draw=black!70, thick},
  card/.style={font=\sffamily\tiny\bfseries, fill=white, inner sep=0.8pt, text=algogreen}}
\providecommand{\entity}[4]{%
  \node[entity] (#1) at (#2) {%
    \renewcommand{\arraystretch}{1.05}\setlength{\tabcolsep}{3pt}%
    \begin{tabular}{l}
      \rowcolor{tierblue}\rule{0pt}{8pt}\textbf{\small #3}\\ \hline
      #4
    \end{tabular}};
  \draw[black!70] (#1.south west) rectangle (#1.north east);}
% \relate{from}{to}{card near from}{card near to}[path options]
\providecommand{\relate}[5][]{%
  \draw[rel, #1] (#2) -- node[card, pos=0.12] {#4} node[card, pos=0.88] {#5} (#3);}
\providecommand{\PK}{\textbf{PK}~}
\providecommand{\FK}{\textit{FK}~}

\begin{tikzpicture}
  % Row 1
  \entity{audit}{-4.2,0}{audit\_logs}{\PK id\\ \FK user\_id\\ action, tool\_name\\ outcome, created\_at}
  \entity{users}{0,0}{users}{\PK id\\ email (unique)\\ password\_hash\\ role}
  \entity{documents}{4.4,0}{documents}{\PK id\\ \FK uploaded\_by\\ title, category, version\\ status, is\_synthetic\\ \FK superseded\_by}
  \entity{chunks}{9,0}{document\_chunks}{\PK id\\ \FK document\_id\\ section, page\\ content\\ embedding vector(384)\\ content\_tsv tsvector}
  % Row 2
  \entity{sess}{-4.2,-4.4}{sessions}{\PK id (random token)\\ \FK user\_id\\ expires\_at}
  \entity{conversations}{0,-4.4}{conversations}{\PK id\\ \FK user\_id\\ title, created\_at}
  \entity{jobs}{4.4,-4.4}{ingestion\_jobs}{\PK id\\ \FK document\_id\\ state, attempts, error\\ created\_at, updated\_at}
  \entity{notices}{9,-4.4}{notices}{\PK id\\ \FK published\_by\\ title, body\\ audience, published\_at}
  % Row 3
  \entity{messages}{0,-8.4}{messages}{\PK id\\ \FK conversation\_id\\ role, content\\ sources (jsonb)}

  \relate{audit.east}{users.west}{N}{1}
  \relate{documents.west}{users.east}{N}{1}
  \relate{chunks.west}{documents.east}{N}{1}
  \relate{sess.north east}{users.south west}{N}{1}
  \relate{conversations.north}{users.south}{N}{1}
  \relate{jobs.north}{documents.south}{N}{1}
  \relate{notices.north west}{users.south east}{N}{1}
  \relate{messages.north}{conversations.south}{N}{1}
\end{tikzpicture}
}
\caption{ER diagram: authentication, documents and chat.}
\label{fig:er-platform}
\end{figure}

\paragraph{Design decisions.}
\begin{itemize}[itemsep=1pt]
  \item The schema is in third normal form. Derived values such as attendance percentage
        are computed by a view (\texttt{attendance\_summary}) instead of being stored, so
        they can never disagree with the underlying records.
  \item Attendance is stored per class session, which allows the calculator to answer
        ``how many more classes'' exactly and lets the generator create realistic patterns.
  \item Integrity is enforced by the database: foreign keys, \texttt{UNIQUE} constraints
        (one enrollment per student per offering), and \texttt{CHECK} constraints (amounts
        $\ge 0$, marks $\le$ maximum).
  \item Indexes: B-tree indexes on every foreign key and on \texttt{(student\_id)} filters;
        an HNSW index~\cite{malkov2020hnsw} on \texttt{embedding}; a GIN index on
        \texttt{content\_tsv} for full-text search.
  \item Every generated row carries \texttt{is\_synthetic = true}, as required by FR-02.
\end{itemize}

Table~\ref{tab:sizes} lists the planned dataset size for 5,000 students.

\begin{table}[H]
\centering
\caption{Planned table sizes for 5,000 synthetic students (approximate)}
\label{tab:sizes}
\begin{tabular}{@{}lrlr@{}}
\toprule
\textbf{Table} & \textbf{Rows} & \textbf{Table} & \textbf{Rows}\\
\midrule
students        & 5,000     & class\_sessions & 18,000\\
faculty         & 150       & attendance      & 1,350,000\\
courses         & 180       & assignments     & 1,600\\
course\_offerings & 400     & submissions     & 110,000\\
enrollments     & 30,000    & exams           & 800\\
invoices        & 10,000    & payments        & 9,000\\
\bottomrule
\end{tabular}
\end{table}

% ===========================================================================
\chapter{Inter-Process Communication}
\label{ch:ipc}

\section{Processes in the System}
At run time CampusGPT consists of five cooperating operating-system processes on the
server side of the browser (Figure~\ref{fig:proc}): the Next.js server, the FastAPI server,
the ingestion worker, the PostgreSQL server (which itself forks one backend process per
connection), and the remote LLM service. They share no memory; all cooperation happens
through the channels numbered \circled{1}--\circled{7}.

\begin{figure}[H]
\centering
\resizebox{\textwidth}{!}{% Process view: OS processes and the IPC channels between them
\begin{tikzpicture}[
  font=\sffamily\scriptsize,
  proc/.style={draw=black!70, thick, rounded corners=3pt, align=center,
               text width=3.5cm, minimum height=1.35cm, inner sep=3pt},
  arr/.style={-{Latex[length=2mm]}, thick, draw=black!75},
  biarr/.style={{Latex[length=2mm]}-{Latex[length=2mm]}, thick, draw=black!75},
  lab/.style={font=\sffamily\tiny, fill=white, inner sep=1.5pt, align=center}]

  \node[proc, fill=tierblue] (browser) at (0,3.6)
    {\textbf{Web Browser}\\student / admin UI\\(JavaScript runtime)};
  \node[proc, fill=tierblue] (next) at (0,7.2)
    {\textbf{Next.js Server}\\Node.js process\\serves pages \& bundles};
  \node[proc, fill=tiergreen] (api) at (6.4,3.6)
    {\textbf{FastAPI Server}\\Uvicorn, Python process\\async event loop};
  \node[proc, fill=tieramber] (worker) at (12.8,3.6)
    {\textbf{Ingestion Worker}\\separate Python process\\one job at a time};
  \node[proc, fill=tiergray] (pg) at (9.6,-0.6)
    {\textbf{PostgreSQL Server}\\postmaster + one backend\\process per connection};
  \node[proc, fill=tiergray, dashed] (llm) at (9.6,7.2)
    {\textbf{LLM Provider}\\remote HTTPS service};

  % (6) pages
  \draw[biarr] (browser.north) -- node[lab] {\circled{6} HTTP\\pages, JS} (next.south);
  % (1) REST + SSE
  \draw[arr] ([yshift=5pt]browser.east) -- node[lab, above] {\circled{1} REST request (JSON)}
    ([yshift=5pt]api.west);
  \draw[arr] ([yshift=-8pt]api.west) -- node[lab, below] {\circled{1} response / SSE stream}
    ([yshift=-8pt]browser.east);
  % (4) LLM
  \draw[biarr] (api.north east) -- node[lab, sloped, above] {\circled{4} HTTPS JSON} (llm.south west);
  % (2) SQL connections
  \draw[biarr] (api.south) |- node[lab, pos=0.3] {\circled{2} SQL\\connection pool} ([yshift=6pt]pg.west);
  \draw[biarr] (worker.south) |- node[lab, pos=0.3] {\circled{2} SQL\\connection} ([yshift=6pt]pg.east);
  % (3) NOTIFY / LISTEN
  \draw[arr, dashed, draw=algogreen] ([xshift=-12pt]api.south east) to[out=-50, in=165]
    node[lab, pos=0.45, text=algogreen] {\circled{3} NOTIFY\\ingest\_jobs} ([yshift=-8pt]pg.north west);
  \draw[arr, dashed, draw=algogreen] ([yshift=-8pt]pg.north east) to[out=15, in=-130]
    node[lab, pos=0.55, text=algogreen] {\circled{3} LISTEN\\wake-up} ([xshift=12pt]worker.south west);
  % (5) shared job table
  \node[lab, draw=black!40, rounded corners=2pt, text width=3.4cm, below=0.25cm of pg]
    {\circled{5} shared \texttt{ingestion\_jobs} table:\\claimed with \texttt{FOR UPDATE SKIP LOCKED}};
\end{tikzpicture}
}
\caption{Process view with inter-process communication channels.}
\label{fig:proc}
\end{figure}

\section{IPC Mechanisms}
Table~\ref{tab:ipc} lists every mechanism, the processes it connects, its communication
pattern, and the reason for choosing it~\cite{stevens1999ipc}.

\begin{longtable}{@{}>{\raggedright\arraybackslash}p{0.6cm}>{\raggedright\arraybackslash}p{3.1cm}>{\raggedright\arraybackslash}p{2.6cm}>{\raggedright\arraybackslash}p{2.6cm}>{\raggedright\arraybackslash}p{5.3cm}@{}}
\caption{IPC mechanisms used in CampusGPT}\label{tab:ipc}\\
\toprule
\textbf{\#} & \textbf{Mechanism} & \textbf{Between} & \textbf{Pattern} & \textbf{Why}\\
\midrule
\endfirsthead
\toprule
\textbf{\#} & \textbf{Mechanism} & \textbf{Between} & \textbf{Pattern} & \textbf{Why}\\
\midrule
\endhead
\bottomrule
\endfoot
\circled{1} & HTTP REST with JSON~\cite{fielding2000rest}, plus Server-Sent
Events~\cite{whatwg-sse} & Browser $\leftrightarrow$ FastAPI & Request--response; server
push stream & Stateless and universal. SSE streams the answer as it is generated, which
lowers perceived latency without the complexity of WebSockets.\\
\circled{2} & PostgreSQL wire protocol over TCP or a Unix domain socket & FastAPI, worker
$\leftrightarrow$ PostgreSQL & Request--response through a connection pool & Reusing
connections avoids a new server process per query; parameterised queries prevent SQL
injection.\\
\circled{3} & \texttt{LISTEN}/\allowbreak\texttt{NOTIFY} \cite{pg-notify} on channel
\texttt{ingest\_jobs} & FastAPI $\rightarrow$ PostgreSQL $\rightarrow$ worker & Asynchronous
publish--subscribe & Wakes the worker instantly without busy polling. A notification sent
inside a transaction is delivered only if the transaction commits.\\
\circled{4} & HTTPS with JSON & FastAPI $\rightarrow$ LLM provider & Remote request--response
with timeout and retry & The only channel that leaves the machine; carries the question,
tool definitions and retrieved text, never raw database access.\\
\circled{5} & Shared \texttt{ingestion\_jobs} table with row locks
(\texttt{FOR UPDATE SKIP LOCKED})~\cite{pg-select} & FastAPI, worker & Shared state with
mutual exclusion & A durable queue: jobs survive crashes, and each job is claimed by exactly
one worker even if several run.\\
\circled{6} & HTTP & Browser $\leftrightarrow$ Next.js server & Request--response & Delivers
pages and JavaScript bundles.\\
\circled{7} & POSIX signals (\texttt{SIGTERM}, \texttt{SIGINT}) & Operating system / process
manager $\rightarrow$ FastAPI, worker & Asynchronous notification & Graceful shutdown: the
worker finishes the current job and the server drains open requests before exiting.\\
\end{longtable}

\paragraph{Why NOTIFY together with a table.} PostgreSQL does not store notifications for
processes that are not listening, so \texttt{NOTIFY} alone would lose jobs uploaded while
the worker is stopped. The table is therefore the durable source of truth, and the
notification is only a wake-up signal. As a safety net, the worker also checks the table
every 30 seconds.

\paragraph{Why not a separate message broker.} Redis, RabbitMQ or Celery would add another
server process and its memory to an 8\,GB development laptop, while PostgreSQL already
provides durable storage, row locking and publish--subscribe.

\section{Worked Example: Document Ingestion}
Figure~\ref{fig:jobs} shows the life cycle of an ingestion job, and the steps below follow
one upload across process boundaries.

\begin{figure}[H]
\centering
\resizebox{0.85\textwidth}{!}{% State machine for an ingestion job
\begin{tikzpicture}[
  font=\sffamily\scriptsize,
  st/.style={draw=black!70, thick, rounded corners=8pt, fill=tierblue, minimum width=2.2cm,
             minimum height=0.85cm, align=center},
  fin/.style={st, fill=tiergreen},
  bad/.style={st, fill=red!8},
  arr/.style={-{Latex[length=2mm]}, thick, draw=black!75},
  lab/.style={font=\sffamily\tiny, fill=white, inner sep=1pt, align=center}]
  \node[st]  (q) at (0,0)    {\textbf{queued}};
  \node[st]  (p) at (4.2,0)  {\textbf{processing}};
  \node[fin] (d) at (8.4,0.9) {\textbf{done}};
  \node[bad] (f) at (8.4,-0.9) {\textbf{failed}};
  \draw[arr] (-1.9,0) -- node[lab, above] {upload} (q.west);
  \draw[arr] (q) -- node[lab, above] {worker claims\\(SKIP LOCKED)} (p);
  \draw[arr] (p.east) -- node[lab, sloped, above, pos=0.4] {chunks stored} (d.west);
  \draw[arr] (p.east) -- node[lab, sloped, below, pos=0.4] {error, attempts = 3} (f.west);
  \draw[arr] (p.south) to[out=-120, in=-60] node[lab, below] {error, attempts $<$ 3: retry} (q.south);
\end{tikzpicture}
}
\caption{State machine of an ingestion job.}
\label{fig:jobs}
\end{figure}

\begin{algobox}{Ingestion across processes}
\begin{enumerate}[leftmargin=*,itemsep=1pt]
  \item \textbf{Browser $\rightarrow$ FastAPI \circled{1}:} the administrator uploads a PDF
        with \texttt{POST /admin/documents}.
  \item \textbf{FastAPI $\rightarrow$ PostgreSQL \circled{2}:} in one transaction, insert the
        \texttt{documents} row and a \texttt{queued} job, then execute
        \texttt{NOTIFY ingest\_jobs} \circled{3}. The API replies \texttt{202 Accepted}
        immediately.
  \item \textbf{PostgreSQL $\rightarrow$ worker \circled{3}:} on commit, the listening worker
        wakes up.
  \item \textbf{Worker $\rightarrow$ PostgreSQL \circled{5}:} claim the oldest queued job with
        \texttt{FOR UPDATE SKIP LOCKED} and set it to \texttt{processing}.
  \item \textbf{Worker (local):} extract text, clean it, split it into chunks and embed them.
  \item \textbf{Worker $\rightarrow$ PostgreSQL \circled{2}:} in one transaction, insert all
        chunks, mark the previous document version as superseded, and set the job to
        \texttt{done}.
  \item \textbf{Browser $\rightarrow$ FastAPI \circled{1}:} the admin console polls
        \texttt{GET /admin/documents/\{id\}} every two seconds and shows the new status.
\end{enumerate}
\end{algobox}

\noindent\begin{minipage}{\linewidth}
\begin{lstlisting}[language=SQL, caption={Enqueueing and claiming a job (shared table and NOTIFY)}, label=lst:queue]
-- FastAPI, inside the upload transaction
INSERT INTO ingestion_jobs (document_id, state) VALUES (:doc_id, 'queued');
SELECT pg_notify('ingest_jobs', :doc_id::text);   -- delivered on COMMIT

-- Worker: claim one job; rows locked by another worker are skipped
UPDATE ingestion_jobs
   SET state = 'processing', attempts = attempts + 1, updated_at = now()
 WHERE id = (SELECT id FROM ingestion_jobs
              WHERE state = 'queued'
              ORDER BY created_at
              LIMIT 1
              FOR UPDATE SKIP LOCKED)
RETURNING id, document_id;
\end{lstlisting}
\end{minipage}


\noindent\begin{minipage}{\linewidth}
\begin{lstlisting}[language=Python, caption={Worker loop: LISTEN with a polling safety net}, label=lst:worker]
async def run_worker(dsn: str) -> None:
    conn = await psycopg.AsyncConnection.connect(dsn, autocommit=True)
    await conn.execute("LISTEN ingest_jobs")
    while not shutdown.is_set():                  # set by the SIGTERM handler
        while (job := await claim_job(conn)) is not None:
            await process(job)                    # extract, chunk, embed, store
        # Sleep until a notification arrives, or 30 s at most
        async for _ in conn.notifies(timeout=30, stop_after=1):
            pass
\end{lstlisting}
\end{minipage}


\section{Streaming Answers to the Browser}
For each question, FastAPI returns a \texttt{text/event-stream} response. The
orchestrator emits typed events: \texttt{status} (for example ``Checking your
attendance\ldots''), \texttt{token} (a piece of the answer), \texttt{sources} (database and
document citations), \texttt{done} and \texttt{error}. Because the browser's
\texttt{EventSource} API supports only \texttt{GET}, the frontend sends a \texttt{POST}
with \texttt{fetch} and reads the event stream from the response body.

\section{Synchronisation and Failure Handling}
\begin{itemize}[itemsep=1pt]
  \item \textbf{Timeouts and retries.} LLM calls time out after 30 seconds and are retried
        at most twice with exponential back-off; afterwards the chat shows a clear error
        while all other pages keep working (FR-16).
  \item \textbf{Worker crash.} A job left in \texttt{processing} for more than ten minutes
        is returned to \texttt{queued} on the next claim, up to three attempts.
  \item \textbf{Atomicity.} Chunks and the job status are written in one transaction, so a
        failure never leaves a half-indexed document visible to search.
  \item \textbf{Back-pressure.} Connection pools cap the number of PostgreSQL backend
        processes; rate limits cap chat requests per user (NFR-S5).
\end{itemize}

% ===========================================================================
\chapter{Programming Concepts}

Table~\ref{tab:concepts} lists the programming concepts used in CampusGPT and where each
appears; selected concepts are illustrated by the listings that follow.

\begin{longtable}{@{}>{\raggedright\arraybackslash}p{3.6cm}>{\raggedright\arraybackslash}p{5.9cm}>{\raggedright\arraybackslash}p{4.9cm}@{}}
\caption{Programming concepts and where they are applied}\label{tab:concepts}\\
\toprule
\textbf{Concept} & \textbf{Where it is used} & \textbf{Purpose}\\
\midrule
\endfirsthead
\toprule
\textbf{Concept} & \textbf{Where it is used} & \textbf{Purpose}\\
\midrule
\endhead
\bottomrule
\endfoot
Abstraction, inheritance, polymorphism & \texttt{LLMProvider} abstract class with Groq,
Gemini and Ollama subclasses (Listing~\ref{lst:provider}) & Orchestrator code is identical
for every provider.\\
Design patterns~\cite{gamma1994} & Adapter (LLM provider, university data source),
Strategy (provider chosen by configuration), Registry (tools), Repository (data access),
Publish--subscribe (\texttt{LISTEN}/\texttt{NOTIFY}) & Known solutions with clear
responsibilities.\\
Decorators and higher-order functions~\cite{pep318} & \texttt{@tool} registration
(Listing~\ref{lst:tools}); FastAPI route decorators & Declarative registration without
boilerplate.\\
Dependency injection & FastAPI \texttt{Depends} for the current user and database session &
Testable handlers; authentication cannot be forgotten.\\
Asynchronous programming~\cite{pep492} & \texttt{async}/\texttt{await} handlers, database
driver and HTTP client & Many concurrent requests on one process while waiting on I/O.\\
Generators and streaming~\cite{pep525} & Async generator producing SSE events
(Listing~\ref{lst:sse}); batch generators in the data generator & Constant memory, output
sent as soon as it exists.\\
Type hints and validation & Pydantic models for requests, responses and tool arguments;
TypeScript on the frontend & Invalid input rejected at the boundary; JSON Schema for tool
definitions generated automatically.\\
Exception handling & Custom hierarchy (\texttt{ToolNotPermitted},
\texttt{ToolValidationError}, \texttt{ProviderUnavailable}) mapped to HTTP status codes &
Consistent, safe error messages.\\
Transactions and ACID & Upload, ingestion and chat logging run as single transactions &
No partial writes.\\
Concurrency control & Row locks with \texttt{SKIP LOCKED}; \texttt{UNIQUE} constraints for
idempotent inserts & Correct behaviour with several workers.\\
Data structures and algorithms & Hash map (tool registry), binary heap (weekly-summary
priority queue), Reciprocal Rank Fusion~\cite{cormack2009rrf}, cosine similarity, sliding-window
chunking & Efficient lookup, ranking and retrieval.\\
Exact arithmetic & \texttt{fractions.Fraction} in the attendance calculator
(Listing~\ref{lst:calc}) & No floating-point rounding errors in eligibility decisions.\\
Secure programming & Argon2id hashing~\cite{biryukov2016argon2}, CSPRNG session tokens
(\texttt{secrets}), parameterised SQL, role checks~\cite{sandhu1996rbac} & Defence against
credential theft, injection and privilege escalation.\\
Text processing & Regular expressions for cleaning extracted PDF text & Removes headers,
footers and broken hyphenation before chunking.\\
Reproducible randomness & Seeded Faker and NumPy generators & Identical synthetic dataset on
every machine.\\
Automated testing & pytest fixtures, parametrised tests, a rolled-back transaction per test
& Fast, isolated, repeatable tests.\\
Frontend components and hooks & React components, \texttt{useState}/\texttt{useEffect},
stream reading with \texttt{ReadableStream} & Responsive UI that renders answers while they
stream.\\
\end{longtable}

\noindent\begin{minipage}{\linewidth}
\begin{lstlisting}[language=Python, caption={Abstraction and polymorphism: the LLM provider adapter}, label=lst:provider]
class LLMProvider(ABC):
    """Interface every model provider implements."""

    @abstractmethod
    async def complete(self, messages: list[Message],
                       tools: list[ToolSpec]) -> LLMReply: ...

class GroqProvider(LLMProvider):
    async def complete(self, messages, tools):
        resp = await self.http.post("/chat/completions",
                                    json=to_openai_format(messages, tools))
        resp.raise_for_status()
        return LLMReply.from_openai_format(resp.json())

def get_provider(settings: Settings) -> LLMProvider:      # Strategy by configuration
    providers = {"groq": GroqProvider, "gemini": GeminiProvider,
                 "ollama": OllamaProvider}
    return providers[settings.llm_provider](settings)
\end{lstlisting}
\end{minipage}


\noindent\begin{minipage}{\linewidth}
\begin{lstlisting}[language=Python, caption={Decorator-based tool registry with server-side permission checks}, label=lst:tools]
TOOLS: dict[str, Tool] = {}

def tool(name: str, roles: set[str], args: type[BaseModel]):
    def register(fn):
        TOOLS[name] = Tool(name=name, roles=roles, args=args, fn=fn)
        return fn
    return register

class AttendanceArgs(BaseModel):
    course: str | None = None

@tool("get_my_attendance", roles={"student"}, args=AttendanceArgs)
async def get_my_attendance(db: AsyncSession, student_id: int, a: AttendanceArgs):
    return await attendance_repo.summary(db, student_id, course=a.course)

async def execute(call: ToolCall, user: User, db: AsyncSession):
    spec = TOOLS.get(call.name)
    if spec is None or user.role not in spec.roles:
        raise ToolNotPermitted(call.name)            # the model cannot bypass this
    args = spec.args.model_validate(call.arguments)  # schema validation
    return await spec.fn(db, user.student_id, args)  # identity from the session
\end{lstlisting}
\end{minipage}


\noindent\begin{minipage}{\linewidth}
\begin{lstlisting}[language=Python, caption={Dependency injection and an async generator streaming SSE events}, label=lst:sse]
@router.post("/chat")
async def chat(req: ChatRequest,
               user: User = Depends(current_user),     # 401 if not logged in
               db: AsyncSession = Depends(get_db)):
    async def events():
        async for event in orchestrator.run(req.question, user, db):
            yield {"event": event.kind, "data": event.model_dump_json()}
    return EventSourceResponse(events())
\end{lstlisting}
\end{minipage}


\noindent\begin{minipage}{\linewidth}
\begin{lstlisting}[language=Python, caption={Deterministic attendance calculation with exact arithmetic}, label=lst:calc]
def classes_needed(attended: int, held: int,
                   target: Fraction = Fraction(3, 4)) -> int:
    """Smallest x such that (attended + x) / (held + x) >= target."""
    if held == 0 or Fraction(attended, held) >= target:
        return 0
    return math.ceil((target * held - attended) / (1 - target))

assert classes_needed(41, 60) == 16      # 57 / 76 = 75%
\end{lstlisting}
\end{minipage}


% ===========================================================================
\chapter{Libraries and Frameworks}

Table~\ref{tab:libs} lists every library CampusGPT depends on. All are free and open source;
exact versions are pinned in \texttt{requirements.txt} and \texttt{package-lock.json}.

\begin{longtable}{@{}>{\raggedright\arraybackslash}p{2.9cm}>{\raggedright\arraybackslash}p{4.4cm}>{\raggedright\arraybackslash}p{2.6cm}>{\raggedright\arraybackslash}p{4.5cm}@{}}
\caption{Libraries, their purpose, and the reason for each choice}\label{tab:libs}\\
\toprule
\textbf{Library} & \textbf{Purpose in CampusGPT} & \textbf{Alternative} & \textbf{Reason for
choice}\\
\midrule
\endfirsthead
\toprule
\textbf{Library} & \textbf{Purpose in CampusGPT} & \textbf{Alternative} & \textbf{Reason for
choice}\\
\midrule
\endhead
\bottomrule
\endfoot
\multicolumn{4}{@{}l}{\textit{Backend}}\\
FastAPI~\cite{fastapi} & HTTP API, routing, dependency injection, OpenAPI documentation &
Flask, Django & Asynchronous by design; validation and API docs built in.\\
Uvicorn & ASGI server that runs the FastAPI process & Gunicorn (WSGI) & Standard
asynchronous server for FastAPI.\\
Pydantic~\cite{pydantic} & Schemas for requests, responses and tool arguments &
marshmallow & Integrated with FastAPI; exports JSON Schema used as LLM tool definitions.\\
SQLAlchemy~2 \cite{sqlalchemy} & ORM and query builder & Raw SQL strings & Typed, composable
queries with bound parameters.\\
psycopg~3~\cite{psycopg} & PostgreSQL driver: async queries, \texttt{LISTEN}/\texttt{NOTIFY},
\texttt{COPY} & asyncpg & Supports all three IPC features we need and is SQLAlchemy's
standard dialect.\\
Alembic & Versioned database migrations & Hand-written SQL scripts & Repeatable,
reversible schema changes.\\
argon2-cffi & Password hashing with Argon2id & bcrypt & Memory-hard algorithm recommended
by OWASP~\cite{owasp-passwords}.\\
httpx & Async HTTP client for LLM providers & requests, vendor SDKs & Non-blocking; one
client for all providers avoids vendor lock-in.\\
sse-starlette & Server-Sent Events responses & WebSockets & One-way streaming is enough and
simpler to secure.\\
slowapi & Rate limiting per user and IP & Custom middleware & Decorator-based limits per
route.\\
\multicolumn{4}{@{}l}{\textit{AI and document processing}}\\
sentence-\allowbreak transformers \cite{reimers2019sbert} & Runs the BGE-small embedding
model~\cite{xiao2024cpack} on the CPU & Paid embedding APIs & Free, offline, and no document
text leaves the machine.\\
pypdf & PDF text extraction & PyMuPDF, pdfplumber & Pure Python with a permissive licence
(PyMuPDF is AGPL).\\
\multicolumn{4}{@{}l}{\textit{Synthetic data generation}}\\
Faker (\texttt{en\_IN}) & Indian names, phone numbers and addresses & Hand-written name
lists & Locale-aware, seedable, large variety.\\
NumPy & Statistical distributions for attendance, marks and fee payments & Python
\texttt{random} & Fast vectorised sampling with a seeded generator.\\
\multicolumn{4}{@{}l}{\textit{Testing and quality}}\\
pytest & Unit, integration and security tests & unittest & Fixtures and parametrisation.\\
Ruff & Linting and formatting & flake8 + black & One fast tool for both.\\
\multicolumn{4}{@{}l}{\textit{Frontend}}\\
Next.js~\cite{nextjs} & React framework: routing and production builds & Vite + React &
File-based routing and built-in optimisation.\\
React~\cite{react} & Component-based user interface & Vue & Largest ecosystem; required by
Next.js.\\
TypeScript & Static typing for frontend code & JavaScript & Catches type errors at build
time.\\
Tailwind CSS & Styling with a consistent design scale & Bootstrap & Small output, consistent
spacing and colour tokens.\\
react-markdown & Renders formatted chat answers & Raw HTML insertion & Does not render raw
HTML, preventing cross-site scripting.\\
Recharts & Attendance and fee charts on the dashboard & Chart.js & Native React
components.\\
\multicolumn{4}{@{}l}{\textit{Infrastructure}}\\
PostgreSQL~16 \cite{postgresql} & Relational database, full-text search, job queue,
notifications & MySQL, MongoDB & One system provides integrity, search and IPC.\\
pgvector~\cite{pgvector} & Vector column type and HNSW index & Pinecone, ChromaDB, FAISS &
No separate vector database; vectors share transactions with their metadata.\\
\end{longtable}

% ===========================================================================
\renewcommand{\bibname}{References}
\begin{thebibliography}{99}
\addcontentsline{toc}{chapter}{References}

\bibitem{lewis2020rag}
P.~Lewis \emph{et al.}, ``Retrieval-augmented generation for knowledge-intensive NLP
tasks,'' in \emph{Advances in Neural Information Processing Systems (NeurIPS 33)}, 2020,
pp.~9459--9474.

\bibitem{pgvector}
A.~Kane, ``pgvector: Open-source vector similarity search for Postgres,'' 2024. [Online].
Available: \url{https://github.com/pgvector/pgvector}

\bibitem{demarco1978}
T.~DeMarco, \emph{Structured Analysis and System Specification}. New York, NY, USA: Yourdon
Press, 1978.

\bibitem{codd1970}
E.~F. Codd, ``A relational model of data for large shared data banks,''
\emph{Communications of the ACM}, vol.~13, no.~6, pp.~377--387, 1970.

\bibitem{chen1976}
P.~P.-S. Chen, ``The entity-relationship model---toward a unified view of data,''
\emph{ACM Transactions on Database Systems}, vol.~1, no.~1, pp.~9--36, 1976.

\bibitem{malkov2020hnsw}
Y.~A. Malkov and D.~A. Yashunin, ``Efficient and robust approximate nearest neighbor search
using hierarchical navigable small world graphs,'' \emph{IEEE Trans. Pattern Analysis and
Machine Intelligence}, vol.~42, no.~4, pp.~824--836, 2020.

\bibitem{stevens1999ipc}
W.~R. Stevens, \emph{UNIX Network Programming, Volume 2: Interprocess Communications},
2nd~ed. Upper Saddle River, NJ, USA: Prentice Hall, 1999.

\bibitem{fielding2000rest}
R.~T. Fielding, ``Architectural styles and the design of network-based software
architectures,'' Ph.D. dissertation, University of California, Irvine, 2000.

\bibitem{whatwg-sse}
WHATWG, ``HTML Living Standard: Server-sent events.'' [Online]. Available:
\url{https://html.spec.whatwg.org/multipage/server-sent-events.html}

\bibitem{pg-notify}
The PostgreSQL Global Development Group, ``NOTIFY,'' \emph{PostgreSQL 16 Documentation}.
[Online]. Available: \url{https://www.postgresql.org/docs/16/sql-notify.html}

\bibitem{pg-select}
The PostgreSQL Global Development Group, ``SELECT: The locking clause,'' \emph{PostgreSQL 16
Documentation}. [Online]. Available:
\url{https://www.postgresql.org/docs/16/sql-select.html}

\bibitem{gamma1994}
E.~Gamma, R.~Helm, R.~Johnson, and J.~Vlissides, \emph{Design Patterns: Elements of Reusable
Object-Oriented Software}. Reading, MA, USA: Addison-Wesley, 1994.

\bibitem{pep318}
K.~D. Smith \emph{et al.}, ``PEP 318 -- Decorators for functions and methods,'' 2003.
[Online]. Available: \url{https://peps.python.org/pep-0318/}

\bibitem{pep492}
Y.~Selivanov, ``PEP 492 -- Coroutines with async and await syntax,'' 2015. [Online].
Available: \url{https://peps.python.org/pep-0492/}

\bibitem{pep525}
Y.~Selivanov, ``PEP 525 -- Asynchronous generators,'' 2016. [Online]. Available:
\url{https://peps.python.org/pep-0525/}

\bibitem{cormack2009rrf}
G.~V. Cormack, C.~L.~A. Clarke, and S.~Buettcher, ``Reciprocal rank fusion outperforms
Condorcet and individual rank learning methods,'' in \emph{Proc. 32nd ACM SIGIR}, 2009,
pp.~758--759.

\bibitem{biryukov2016argon2}
A.~Biryukov, D.~Dinu, and D.~Khovratovich, ``Argon2: New generation of memory-hard functions
for password hashing and other applications,'' in \emph{Proc. IEEE European Symposium on
Security and Privacy (EuroS\&P)}, 2016, pp.~292--302.

\bibitem{sandhu1996rbac}
R.~S. Sandhu, E.~J. Coyne, H.~L. Feinstein, and C.~E. Youman, ``Role-based access control
models,'' \emph{Computer}, vol.~29, no.~2, pp.~38--47, 1996.

\bibitem{fastapi}
S.~Ram\'irez, ``FastAPI documentation.'' [Online]. Available:
\url{https://fastapi.tiangolo.com/}

\bibitem{pydantic}
Pydantic, ``Pydantic documentation.'' [Online]. Available:
\url{https://docs.pydantic.dev/}

\bibitem{sqlalchemy}
SQLAlchemy, ``SQLAlchemy 2.0 documentation.'' [Online]. Available:
\url{https://docs.sqlalchemy.org/en/20/}

\bibitem{psycopg}
D.~Varrazzo, ``Psycopg 3 documentation.'' [Online]. Available:
\url{https://www.psycopg.org/psycopg3/docs/}

\bibitem{owasp-passwords}
OWASP Foundation, ``Password Storage Cheat Sheet.'' [Online]. Available:
\url{https://cheatsheetseries.owasp.org/cheatsheets/Password_Storage_Cheat_Sheet.html}

\bibitem{reimers2019sbert}
N.~Reimers and I.~Gurevych, ``Sentence-BERT: Sentence embeddings using Siamese
BERT-networks,'' in \emph{Proc. EMNLP-IJCNLP}, 2019, pp.~3982--3992.

\bibitem{xiao2024cpack}
S.~Xiao \emph{et al.}, ``C-Pack: Packed resources for general Chinese embeddings,'' in
\emph{Proc. 47th ACM SIGIR}, 2024, pp.~641--649.

\bibitem{nextjs}
Vercel, ``Next.js documentation.'' [Online]. Available: \url{https://nextjs.org/docs}

\bibitem{react}
Meta Open Source, ``React documentation.'' [Online]. Available: \url{https://react.dev/}

\bibitem{postgresql}
The PostgreSQL Global Development Group, ``PostgreSQL 16 documentation.'' [Online].
Available: \url{https://www.postgresql.org/docs/16/}

\end{thebibliography}

\end{document}
